The proof is in Gauge Theory Gravity: for a vacuum field the curvature operator is its Weyl part, $\mathcal{R} = \mathcal{W}$, a linear map on physical bivectors that is symmetric, $\langle A\,\mathcal{W}(B)\rangle_{0} = \langle\mathcal{W}(A)\,B\rangle_{0}$, and trace-free, $\partial_{a}\cdot\mathcal{W}(a \wedge b) = 0$ (Fact \claimref{synthesis.imported-gtg-dictionary}, clauses (iv) and (v)); the null tetrad and the Weyl scalars are those of Fact \claimref{pairing.imported-weyl-scalars-petrov-type-schwarzschild}, $\Psi_{0} = C(\ell, m, \ell, m)$, $\Psi_{1} = C(\ell, n, \ell, m)$, $\Psi_{2} = C(\ell, m, \mb, n)$, with the flat product of physical vectors extended bilinearly over the complex unit of $m$.

\emph{$C$ as a bilinear form of the Weyl map.} For a bivector $B$ and vectors $w, z$, $(B\cdot w)\cdot z = B\cdot(w \wedge z) := \langle B\,(w\wedge z)\rangle_{0}$ (for $B = a \wedge b$ both sides are $(a\cdot z)(b\cdot w) - (a\cdot w)(b\cdot z)$). Hence
\[
C(u, v, w, z) = R(u, v)\cdot(w \wedge z) = -\,\mathcal{W}(u \wedge v)\cdot(w \wedge z).
\]
So $C$ is antisymmetric in $(u, v)$ and in $(w, z)$, and the symmetry of $\mathcal{W}$ gives the pair symmetry $C(u, v, w, z) = C(w, z, u, v)$.

\emph{The trace in the null tetrad.} For any frame $\{f_{a}\}$ with reciprocal frame $\{f^{a}\}$, $\partial_{a} = \sum_{a} f^{a}\,(f_{a}\cdot\partial)$, and for a vector $c$ and bivector $B$, $(c\cdot B)\cdot v = -B\cdot(c \wedge v)$. So
\[
0 \;=\; v\cdot\bigl(\partial_{a}\cdot\mathcal{W}(a \wedge u)\bigr) \;=\; \sum_{a} C(f_{a}, u, f^{a}, v) .
\]
For the tetrad $\ell, n, m, \mb$ with $\ell\cdot n = 2$, $m\cdot\mb = -2$ and the other products zero, the reciprocal frame is $\ell \mapsto n/2$, $n \mapsto \ell/2$, $m \mapsto -\mb/2$, $\mb \mapsto -m/2$, and the trace condition reads
\[
\tfrac12\bigl[C(\ell, u, n, v) + C(n, u, \ell, v)\bigr] - \tfrac12\bigl[C(m, u, \mb, v) + C(\mb, u, m, v)\bigr] = 0 \qquad (\mathrm{T}_{u,v})
\]
for all $u, v$ (for another normalisation of the tetrad the coefficients scale accordingly).

\emph{The six classes.} The class of the loop $(u, v)$ is $C(u, v, \ell, m)$. The loops $(\ell, m)$ and $(\ell, n)$ give $\Psi_{0}$ and $\Psi_{1}$ by definition. The loop $(n, \mb)$ gives $C(n, \mb, \ell, m) = C(\ell, m, n, \mb) = -C(\ell, m, \mb, n) = -\Psi_{2}$, by the pair symmetry and the antisymmetry in the last pair. In $(\mathrm{T}_{\ell, \ell})$ the two $\ell$-terms vanish by antisymmetry, so $C(m, \ell, \mb, \ell) = -C(\mb, \ell, m, \ell)$, while the pair symmetry makes them equal; hence both vanish, and $C(\ell, \mb, \ell, m) = C(\ell, m, \ell, \mb) = C(m, \ell, \mb, \ell) = 0$. $(\mathrm{T}_{m, \ell})$ leaves $C(\ell, m, n, \ell) = C(\mb, m, m, \ell)$; the left side is $-C(\ell, m, \ell, n) = -C(\ell, n, \ell, m) = -\Psi_{1}$ and the right side is $C(m, \mb, \ell, m)$, so the loop $(m, \mb)$ has class $-\Psi_{1}$. $(\mathrm{T}_{m, m})$ leaves $C(\ell, m, n, m) + C(n, m, \ell, m) = 0$, and the two terms are equal by the pair symmetry, so the loop $(n, m)$ has class $0$.

\emph{The equivalence.} The loops in the null cone of $\ell$ are the combinations of $(\ell, m)$, $(\ell, \mb)$ and $(m, \mb)$, with classes $\Psi_{0}$, $0$, $-\Psi_{1}$; they all vanish exactly when $\Psi_{0} = \Psi_{1} = 0$, which by the Fact is the condition that $\ell$ is a repeated principal null direction. Then the radial loop $(\ell, n)$ has class $\Psi_{1} = 0$, and of the six only $(n, \mb)$, with $-\Psi_{2}$, can be nonzero.

The check computes, on all ten basis Weyl maps of $\Cl(1,3)$ and on a generic operator, the identity $(B\cdot w)\cdot z = B\cdot(w \wedge z)$, the reciprocal-frame form of the trace (it equals $v\cdot(\partial_{a}\cdot X(a \wedge u))$ for every bivector map $X$, Weyl or not, and vanishes on the Weyl maps) and the pair symmetry, and then the six identities, both in the adapted tetrad and in the unit-normalised one (check\_gtg\_trace\_and\_pair\_symmetry); the earlier functions verify the table on the Weyl space and on the Schwarzschild operator.

\emph{In components.} The same argument in an orthonormal frame is the source's: the trace is $\eta^{ab}C(e_{a}, u, e_{b}, v) = 0$, written in the null tetrad as $(\mathrm{T}_{u,v})$ up to the normalisation.
